% =====================================================================
% CHAPTER 6: Thesis Management and Finance (9-chapter RUET MTE template)
% Assembled per thesis/RESTRUCTURE_PLAN.md from:
%   - old appendices/app_b.tex "Project Timeline" (phases as a table;
%     the Gantt figure fig:gantt itself moved to Chapter 1, Section 1.5.1
%     and is only referenced here, never duplicated)
%   - old appendices/app_b.tex "Expenses" (tab:expenses, moved intact)
%   - old appendices/app_b.tex "Economic Analysis" (split into hardware
%     vs cloud/software budget)
% New prose: 6.1 Introduction, 6.5 Summary, connective text.
% PO demonstrated: PO11 (project management and finance).
% =====================================================================
\chapter{Thesis Management and Finance}

\section{Introduction}
This chapter records the schedule, allocation of resources, and
expenditure associated with the thesis. The main project period extended
from December 2025 to July 2026, building on the paper collection and
literature review undertaken during the fourth-year odd semester (the
seventh semester).
Model training and experimental studies used rented cloud GPUs, while a
stereo camera and a Jetson Orin Nano were obtained for data acquisition
and deployment testing. Section 6.2 describes the management of the work
and its seven phases, Section 6.3 presents the overall expenditure, and
Section 6.4 separates the hardware, cloud-computing, and software costs.

\FloatBarrier
\section{Management of the Thesis}
The work was divided into phases with defined outputs and deliberate
overlap where activities could proceed in parallel. The earlier literature
review informed the selection of candidate models; setup and preliminary
training then established the experimental pipeline. Real-data capture
continued while the architecture was refined through ablation studies.
The selected model subsequently passed through benchmarking,
three-dimensional reconstruction, and final evaluation before the results
were consolidated in the thesis. This arrangement allowed data
preparation, experimentation, and writing to advance without weakening
the dependency between design decisions and experimental evidence.

\subsection{Time Frame}
The main project schedule ran from December 2025 to July 2026. Paper
collection and literature review had already started in the fourth-year
odd semester (the seventh semester), before the scheduled
model-development work began. Table
\ref{tab:timeframe} lists the periods and principal outputs of the seven
project phases. The corresponding Gantt chart is included in the work
plan of Chapter 1.

\begin{table}[!htbp]
\centering
\caption{\textbf{Time frame of the thesis.} The seven work phases from December 2025 to July 2026.}
\label{tab:timeframe}
\small
\begin{tabular}{p{4.2cm}p{3.0cm}p{6.3cm}}
\toprule
Phase & Period & Output\\
\midrule
1. Model selection and setup & Dec. 2025 to Jan. 2026 & Selected baseline models and a working experimental environment\\
2. Pre-training and fine-tuning & Dec. 2025 to Mar. 2026 & Initial checkpoints and a stable training procedure\\
3. Real-data capture & Jan. to Mar. 2026 & Calibrated and rectified stereo image pairs\\
4. Ablation studies & Feb. to Apr. 2026 & Evidence for the final architectural and training choices\\
5. Benchmarking and 3D reconstruction & Mar. to May 2026 & Accuracy, efficiency, and reconstruction results\\
6. Result analysis and evaluation & Apr. to Jun. 2026 & Consolidated in-domain, zero-shot, robustness, and deployment findings\\
7. Thesis writing and documentation & May to Jul. 2026 & Final thesis book and supporting documentation\\
\bottomrule
\end{tabular}
\end{table}

\FloatBarrier
\section{Overall Budget}
Table \ref{tab:expenses} records the approximate project expenditure.
Cloud computation accounted for about one third of the total: the A100
training workload cost approximately USD 100, and the other GPU instances
used for ablation and evaluation cost approximately USD 40. The stereo
camera and Jetson Orin Nano together cost about BDT 35,000, equivalent to
approximately USD 280 at the conversion used for this project record.

\begin{table}[!htbp]
\centering
\caption{\textbf{Project expenses.} Values are rounded records of the actual development expenditure; the BDT equivalents use approximately BDT 125 per USD.}
\label{tab:expenses}
\small
\begin{tabular}{p{8.4cm}rr}
\toprule
Item & Cost (USD) & Cost (BDT)\\
\midrule
Cloud A100 usage for training & 100 & 12,500\\
Other cloud GPUs for ablation and evaluation & 40 & 5,000\\
Stereo camera and Jetson Orin Nano & 280 & 35,000\\
Software licences & 0 & 0\\
\midrule
Total & 420 & 52,500\\
\bottomrule
\end{tabular}
\end{table}

\FloatBarrier
\section{Component/Software Budget}
The budget of Table \ref{tab:expenses} divides into a hardware component
and a cloud and software component.

The hardware expenditure comprised the stereo camera and the Jetson Orin
Nano used for real-data acquisition and on-device testing. Together these
items cost approximately BDT 35,000 (USD 280). Treating the camera and
edge computer as development hardware, rather than describing the camera
alone as the entire deployed system, gives a more faithful account of the
resources required to validate the proposed pipeline.

The cloud-computing expenditure was approximately USD 140. Of this, about
USD 100 was spent on A100 instances for the principal training workload,
while approximately USD 40 covered the other GPU instances used for
ablation studies, intermediate checks, and final evaluation. These are
usage costs accumulated over the development period, rather than the
price of a single nominal one-day run. Public on-demand GPU rates vary by
accelerator, region, and storage allocation, so the recorded expenditure
is more informative here than a provider's current hourly tariff.

The software environment described in Chapter 7 is based on open-source
tools and incurred no licence fee. Moreover, cloud resources were needed
for development rather than for normal operation: inference is performed
locally on the Jetson device. Consequently, deployment does not introduce
a recurring per-frame cloud-computing charge, although ordinary hardware
maintenance and electrical consumption remain.

\FloatBarrier
\section{Summary}
This chapter summarized the management and finances of the thesis. The
main schedule extended from December 2025 to July 2026, following the
literature work begun in the preceding semester, and comprised seven
overlapping phases from model setup to final documentation. Total
expenditure was approximately USD 420 (BDT 52,500): USD 140 for cloud GPU
usage and USD 280 for the stereo camera and Jetson Orin Nano. The use of
open-source software avoided licence fees, and local inference on the
Jetson device avoids a recurring cloud-inference charge during operation.
